% !TeX root = main.tex
\section{绪论}
\label{sec:introduction}

\subsection{作物育种的表型瓶颈与表型组学}
\label{sec:research-background}

作物育种需要从遗传背景不同的材料中，选出在目标环境下具有优良性状的基因型。随着基因组测序与分子标记技术的发展，研究者已能够在大规模种质群体中系统描述遗传变异，但判断这些变异是否具有农艺价值，仍须以植株的实际表现为依据。以耐热性鉴定为例，一份材料在高温年份保持较高产量，既可能是因为其开花期结实受高温影响较小，也可能是因为开花较早，避开了高温最强的时段。两者对生产都可能有利，但仅凭收获产量，难以判断材料具有哪一种优势，还需结合开花时间、受热条件和结实表现作出解释。同一份种质的基因型记录通常可以在不同试验间复用，而其表型需要在具体环境和发育阶段重新测定。因此，制约育种评价的一个重要因素，是能否在足够多的材料上获得与选育目标相符、彼此可比的性状记录\papercite{furbank2011phenomics}\papercite{araus2018geneticgain}。

表型鉴定的困难，首先体现在群体规模与调查精度之间的取舍。扩大种质群体可以覆盖更丰富的遗传变异，增加生物学重复有助于区分材料间差异与个体波动，开展多年或多点试验则可以检验某一材料的优势能否在不同环境中重现。每增加一个试验地点或观测时期，都意味着要对相应材料再完成一轮调查。对于株高等可以直接量取的性状，主要工作在于逐株测量和记录；对于器官数量、冠层组成和组织受损程度，调查者还需辨认相互遮挡的结构，或依据统一标准估计连续变化的性状。为在有限时间内完成调查而减少重复，会降低材料均值的估计精度；将连续变化的受损程度简化为少数等级，虽然可以加快记录，却可能使表现接近的材料落入同一等级。育种中的表型瓶颈，正体现在群体扩大以后，能够调查的材料数量与每份材料能够获得的性状信息难以同时增加\papercite{araus2018geneticgain}\papercite{watt2020phenotyping}。

对于生长过程中不断变化的性状，调查还受观测窗口的约束。籽粒可以在收获后保存，再按统一流程测定粒重；叶片展开、冠层衰老或胁迫恢复等过程一旦错过，则难以补测。两份材料在处理结束时可能具有相近的枯黄程度，但其中一份较早出现症状、随后缓慢加重，另一份则可能在处理后期才迅速枯黄。一次终点调查足以比较当时的枯黄程度，却不能区分症状出现的早晚与此后的变化速率。观测日期相同也不意味着发育状态相同：在统一日期仍保持绿色的材料，可能衰老较慢，也可能只是较晚进入衰老阶段。因而，评价特定天气过程后整个群体的表现，应尽量缩短同一轮调查的时间跨度；比较材料自身的衰老进程，则需记录关键物候，并在相应发育阶段重复观测。提高测量通量的价值之一，就是使研究者有条件按照性状变化的时间安排调查，而不必主要依据完成一轮人工测量所需的时间安排试验\papercite{jiang2024modernphenomics}\papercite{roth2025fieldphenotyping}。

植物表型组学（plant phenomics）将群体试验、传感观测与定量分析结合起来，系统研究不同基因型在特定环境中的形态结构、生长过程和生理表现\papercite{fiorani2013phenotyping}\papercite{yang2020phenomics}\papercite{jiang2024modernphenomics}。非破坏性成像是其中一种重要的观测方式。传统调查通常将当时的判断记录为一个数值或等级；成像则在同一时刻保留植株轮廓、器官分布和可见组织状态，使田间采集与后续判读能够分开进行。对于短时间内需要调查的大群体，可以先在统一窗口内完成拍摄，再按相同规则处理图像。对同一植株或小区重复成像，还能保持前后观测的对应关系，使研究者可以追踪材料差异何时出现、随后扩大还是缩小。与之相配合，人工计数、收获称量和生理测定为图像指标提供直接参照。两类观测各有用途：影像使同一群体可以反复调查，直接测定则帮助确定影像中可见的变化具有何种农学含义\papercite{roth2025fieldphenotyping}。

成像使一次调查可以保留多项信息，也使测量对象的界定更加重要。以一幅小区俯视图为例，统计植株像素占整幅图像的比例，得到的是给定视角下的冠层投影覆盖度；以植株像素为分母计算绿色组织所占比例，描述的则是可见冠层的组织状态。即使前一指标相同，后一指标也可能随叶片衰老而明显改变。类似地，一幅图像中能够辨认的器官数取决于拍摄视角和遮挡程度，不能无条件地视为植株的器官总数。由图像构建表型时，因而需要说明调查的是单株还是小区、统计的是可见组织还是整株性状，以及多个重复如何汇总为材料水平的观测值。只有这些条件一致，不同材料或时点上的数值才具有明确的比较基础\papercite{fiorani2013phenotyping}\papercite{jiang2024modernphenomics}。

当田间采集能够在短时间内完成时，逐幅判读往往成为随后耗时最多的环节。以器官计数为例，人工在田间逐株辨认改为在屏幕上逐图标记，能够保留核查依据，却未必减少辨认每个目标所需的劳动。自动解析的难点还在于，同一种植物组织在田间并不总具有稳定的外观：叶片可因发育或胁迫而失绿，阴影会改变局部亮度，土壤与残茬又可能呈现接近枯叶的颜色。在背景单一、光照稳定时有效的颜色阈值，换到另一批材料上便可能改变前景的判定范围。深度学习为这类问题提供了利用标注样本建立判读规则的方法。模型可以结合目标自身的颜色、纹理和轮廓，以及它与周围叶片、冠层和背景的空间关系，从而在单一外观特征不足以区分类别时利用更多图像证据。对于高通量表型调查，这种方法的价值在于将同一套明确的标注标准用于批量图像，并保留能够返回原图检查的识别结果\papercite{murphy2024deepphenotyping}\papercite{ubbens2025firsttenyears}。

不同性状需要不同形式的识别结果。测量冠层覆盖，需要逐像素区分植株与背景；估计某类组织在冠层中的比例，还要在已确定的植株范围内区分组织状态；调查离散器官的数量，则可以记录每个可辨识目标的位置，由同一组坐标获得空间分布和计数。输出形式的选择关系到测量误差能否被发现。例如，两种方法即使得到相近的前景面积，其中一种仍可能遗漏冠层外围的细窄叶片，同时将面积相近的背景误判为植株。仅看面积，两类错误可以相互抵消，叠加原图检查边界时才能分辨。器官计数也有类似情形：漏掉一个真实目标，同时在邻近位置多计一个目标，预测总数仍然正确。区域掩膜和点坐标的意义因此不只在于计算一个性状值，还在于保留其空间依据，使研究者能够判断数值上的接近是否来自对同一组织或器官的正确识别\papercite{murphy2024deepphenotyping}。

对于育种材料的比较，图像识别误差是否随株型或组织状态而变化，与误差的平均大小同样值得关注。若一种计数方法在器官稀疏的材料上较准确，却随器官密度增加而持续漏计，群体高值端的差异就可能受到压缩；若分割方法更容易识别绿色组织，而遗漏部分已失绿的植株区域，受损程度较重的材料便可能得到系统性偏低的枯黄比例。这类偏差不会因增加同类图像而自行消失。相应地，图像表型方法需要从两个层面接受评价：在图像层面，依据人工标注检查哪些边界、组织或器官容易识别错误；在材料层面，通过配对调查或重复测量，检验这些错误是否足以影响性状差异的估计。前者为算法改进指明具体对象，后者决定所得性状值适用于怎样的材料比较\papercite{araus2018geneticgain}\papercite{murphy2024deepphenotyping}。

由此，作物育种中的表型瓶颈既体现为有限时间内能够观测多少材料，也体现为所获记录能否转化为定义清楚、误差可评估的性状值。非破坏性成像扩大了群体调查和重复观测的范围，图像分析方法则决定其中哪些组织与器官可以被稳定辨识，以及这些识别结果能够支持何种测量。水稻、小麦和玉米研究已围绕冠层区域提取、器官定位与组织状态评价积累了较丰富的实例，为比较不同视觉任务的标注方式、算法输出和验证依据提供了基础。


\subsection{表型组学在禾本科粮食作物中的研究范式}
\label{sec:cereal-phenomics}

在禾本科粮食作物的表型研究中，同一幅图像可以支持多项测量，但不同性状所需的视觉证据并不相同。以抽穗期的小麦冠层为例，估计覆盖度，需要确定植株与地表背景的边界；调查穗数，需要将相邻麦穗辨认为不同个体；评价花后衰老，则需要区分绿色与衰老组织，并保留各自的投影范围。仅有植株掩膜，无法判断连成一片的区域内包含几个麦穗；仅有器官坐标，也不足以计算枯黄组织所占的面积。由此，水稻、小麦和玉米的图像表型研究逐渐形成了一种共同的研究思路：先依据性状确定观测对象，再选择与之相符的区域、个体或关键点标注，最后分别评价图像识别结果与据此得到的性状值\papercite{murphy2024deepphenotyping}\papercite{wang2025gwfss}\papercite{yang2025cerealdensity}。深度学习方法在这一过程中的作用，是使原本依赖人工逐幅辨认的组织与器官，能够在保留空间位置的条件下接受批量测量。

在冠层与器官形态测量中，像素级语义分割使研究者可以按不同目的组织同一套标注。全球小麦器官语义分割数据集 GWFSS 对叶片、茎秆、麦穗和土壤分别进行标注\papercite{wang2025gwfss}。将叶、茎和穗合并为前景，可以计算作物在视野中的投影覆盖比例；保留各器官类别，则可比较抽穗前后可见冠层的组成变化。这两种测量对同一错误的反应不同：将麦穗像素误分为叶片，不会改变植株前景的总面积，却会低估穗部占比；将伸向冠层外围的窄叶误判为土壤，则会同时改变植株轮廓与覆盖度。水稻多类别语义分割数据集 RiceSEG 为复杂稻田场景提供了相应的像素级参照，使不同方法可以在相同图像与类别定义下接受比较\papercite{zhou2025riceseg}。这类公开标注的意义，不仅在于提供训练样本，也在于使性状测量中的偏差能够追溯至具体区域和错误类别。

禾本科叶片的形态还使冠层分割面临尺度上的困难。株丛中心通常占据较大面积，利用周围结构有助于判断一片区域是否属于植株；叶片向外延伸时逐渐变窄，末端与背景的边界又依赖局部细节。若模型较完整地识别了冠层主体，却漏掉多条外围窄叶，整幅图像的区域重叠指标仍可能较高。反过来，局部颜色或线状边缘本身也不足以确定叶片边界，因为田间残茬可能呈现相近的外观。因而，分割算法是否适合表型测量，既与整幅图像上的识别性能有关，也与误差集中在哪些结构上有关。对于比较叶宽或株型不同的材料，这一区别尤其重要：若窄叶材料持续丢失较多外围组织，即使群体上的平均分割精度较高，材料间投影面积的差异仍可能受到系统性影响。

组织状态测量进一步要求模型在确定前景时保留已经失绿的组织。Serouart 等提出的 SegVeg 先利用 U-Net 从背景中提取植被，再利用多个颜色空间的像素特征，通过支持向量机区分绿色和衰老植被\papercite{serouart2022segveg}。在其研究中，绿色植被比例的估计优于衰老植被比例，衰老组织与背景之间的混淆是较难解决的问题。这一结果说明，颜色分类准确与整幅图像的组织比例准确是两个不同层次的判断。在已正确提取的植株像素上，分类器可以较好地区分绿色和衰老组织；但若部分枯叶此前已被排除在前景之外，它们便不会进入后续统计。相反，与枯叶颜色相近的裸土一旦被纳入植株，也可能造成枯黄面积高估。对于这类方法，分别检查前景边界、组织类别及逐图计算的比例，比仅报告一个总体分类准确率更能说明误差的来源。

水稻穗部研究展示了怎样进一步检验分割结果是否足以支持器官测量。Yang 等开发的 PanicleNeRF 以智能手机采集的田间稻穗视频为输入，结合 YOLOv8 与分割基础模型 Segment Anything Model（SAM）提取不同视角下的穗部区域，再利用神经辐射场重建单穗的三维结构\papercite{yang2024paniclenerf}。研究者分别评价了二维分割边界和重建后的穗长；在其籼稻与粳稻样本中，穗长测量的相对均方根误差分别为 2.94\% 和 1.75\%。这一设计把模型是否找准器官与所得几何性状是否准确区分开来。对于冠层覆盖，二维投影本身就是观测对象；对于姿态弯曲、朝向各异的稻穗，单一视角下的投影长度却会随朝向改变。引入多视角影像的依据，正是目标性状所需的空间信息超出了单幅图像能够稳定提供的范围。

对于叶尖等离散结构，点坐标比区域边界更贴近调查对象。标注一片叶的可见轮廓有明确含义，而标注“叶尖区域”还需额外约定从末端向叶基延伸多远。记录末端位置，可以在减少这类边界分歧的同时，保留图像中的可见叶尖数及其空间分布。基于点标注的算法又有不同输出：密度估计通过对预测密度求和获得数量，直接点预测则输出候选目标坐标，可以逐点检查哪些真实叶尖未被找到、哪些预测落在了相邻叶缘或交叠处\papercite{song2021p2pnet}。这一差别在密集冠层中具有实际意义。一片叶被另一片叶遮挡时，可见叶段也会突然终止；判断它是否为自然叶尖，需要结合周围叶片的走向，而准确标出末端又要求保留足够细的空间信息。因而，计数误差与限定距离内的点匹配结果分别评价数量和位置，不能用其中一项替代另一项。

小麦苗期研究还提供了两条互补的叶尖算法开发路径。一条着眼于训练标注的获取。Li 等利用三维植物模型生成 15 万幅小麦幼苗模拟 RGB 图像及超过 200 万个叶尖标注，并结合域适应方法开展真实图像中的叶尖识别与叶片计数\papercite{li2023selfsupervisedphenotyping}。通过模拟生成图像，器官位置可以随渲染过程一并获得，减少了逐点人工标记的需求；但模拟植株与田间幼苗在叶片姿态、背景和成像条件上的差异，仍需通过真实样本评价。另一条路径着眼于器官识别结果如何转化为农艺调查量。Yang 等利用点预测模型从小麦和大麦的亚毫米级 RGB 影像中获取叶尖密度，再结合积温与出叶进程估计早期植株密度\papercite{yang2025cerealdensity}。新苗出土与已有幼苗展开新叶都会增加叶尖数，因而必须引入发育信息才能由叶尖动态推算株数。前一项研究解决的是如何学习可靠的叶尖识别，后一项则说明，即使叶尖已经识别准确，器官数与最终目标性状之间的关系仍需单独建立。

玉米研究将器官辨识拓展到了无人机观测尺度。Chen 等开发的 DP-MaizeTrack 以苗期无人机影像为输入，在 YOLOv8 基础上增强多尺度特征提取，用于玉米单株和叶片检测，并将检测结果用于区域统计与叶龄分析\papercite{chen2025dpmaizetrack}。与近距离小麦叶尖图像相比，无人机影像一次可以覆盖较多植株，但单片叶所占像素较少，植株和叶片在同一画面中构成两种尺度的目标。该研究在单株检测任务上报告的 \(\mathrm{mAP}_{50}\) 为 94.9\%；这一数值评价的是预测框与植株目标的匹配，不能直接解释为叶片计数准确率。这个例子也说明，发表结果中的“检测精度”只有在明确检测对象和匹配规则后才具有可比性。单株调查、叶片计数与叶尖定位即使取自同一批玉米影像，其参照也应分别落在植株个体、叶片个体和叶片末端上。

对于冠层上方难以看清的器官，改变取像部位有时比更换网络结构更为直接。水稻生长期的俯视图主要记录上层叶片，而区分分蘖所需的茎秆信息集中于株丛基部。Zhou 等开发了基于 Transformer 的水稻分蘖表型模型 TillerPET，选择收获后暴露的田间稻茬作为观测对象，利用茎秆位置测定分蘖数与株丛紧凑度\papercite{zhou2025tillerpet}。定位到的茎秆数量用于估计收获时保留的分蘖数，其空间排列则保留了单纯计数无法描述的株型差异。两份材料可以具有相近的分蘖数，却因茎秆分布疏密不同而呈现不同的紧凑度。与叶尖定位相比，稻茬定位同样输出坐标，但两者需要识别的生物学结构和适宜的取像时期均不同。能否将一种点定位算法迁移至另一器官，首先取决于新图像中是否保留了区分该器官个体的视觉依据。

上述方法从单项试验走向跨场景应用，还须检验训练条件改变后识别性能能否保持。田间图像中的叶色、器官尺度和遮挡程度随材料与生育期而变，土壤背景、拍摄距离与光照条件又可能随地点和批次改变。David 等组织的全球小麦穗部检测挑战采用不同地区采集的田间图像比较检测模型，将新场景中的表现纳入算法评价\papercite{david2023wheatchallenges}。对多来源数据而言，验证设计决定了结果所能支持的判断：在同一来源内部训练与验证，可以比较各算法在该类图像上的表现；合并多个来源后评价，可以考察同一模型处理混合样本的能力；若要判断模型能否用于新的地点或作物，则需保留训练阶段未见的来源作为检验。尤其当各数据集的样本数量和器官密度不同时，相近的汇总误差可能对应不同的分来源表现。公开数据集因此既提供共同的比较基准，也要求研究者明确其结果适用于哪些成像条件。

经过识别与验证的图像指标，还可以沿时间和材料两个维度用于群体研究。Yu 等在两个田间生长季对 565 份小麦材料从开花至成熟开展无人机多光谱观测，以相对持绿评分描述材料间花后绿色冠层维持程度的差异\papercite{yu2024staygreen}。Li 等则在三个环境中将 262 个小麦重组自交系的无人机 RGB 影像与地面旗叶 SPAD 测量相结合，为解释冠层和特定叶位的衰老变化提供了不同尺度的参照\papercite{li2025wheatsenescence}。在水稻和玉米中，图像性状也已用于群体遗传关联分析及跨环境产量预测\papercite{yang2014ricegwas}\papercite{adak2024fieldmaize}\papercite{adak2026temporalmaize}。这些应用回答的是图像性状能否表征材料差异、能否为进一步分析提供信息，而分割边界与器官位置是否准确，仍须回到原图及相应标注上检验。

禾本科粮食作物研究对黑麦草最有价值的参照，是一组可以用相近标注和评价方法共同检验的视觉问题。水稻和小麦与黑麦草均具有细长叶片，算法都需在复杂背景中保留窄叶边缘，判断交叠处的轮廓中断是否为真实叶尖，并在株丛密集时区分相邻目标。因此，公开的禾本科作物掩膜与叶尖点标注适合用于比较冠层分割和器官定位方法；模型能否在黑麦草上保持这些能力，还需依据黑麦草图像中的直接观测加以验证。可以跨作物借鉴的是识别对象的表示方式、算法设计和评价依据，而可见叶尖数、冠层投影和绿色组织比例在牧草育种中分别代表何种性状，则取决于黑麦草自身的株丛结构与利用方式。

\subsection{黑麦草的育种价值、目标性状与表型测定挑战}
\label{sec:ryegrass-breeding-traits}

多年生黑麦草（\textit{Lolium perenne} L.）是温带地区重要的冷季型牧草和草坪草。作为饲用作物，其主要利用对象是一个生长季内可经多次放牧或刈割的地上茎叶。对牧场而言，品种价值不仅取决于全年能够生产多少干物质，还取决于饲草在哪个季节形成、具有怎样的营养品质，以及其中有多少能够被家畜利用。两个全年产量相近的品种，如果一个主要在春季旺盛生长，另一个在夏秋供草不足时仍能保持较高产量，对牧场饲草供应的价值便可能不同。McEvoy 等针对爱尔兰草地畜牧生产体系建立的多年生黑麦草品种经济评价指数，分别考虑春季、生长季中期和秋季的干物质产量，并纳入消化率、青贮刈割产量和草地持久性\papercite{mcevoy2011ryegrasseconomic}。Tubritt 等进一步将放牧利用性纳入品种评价，使品种间可利用程度的差异也能体现在经济评价中\papercite{tubritt2021grazingutilization}。这些指标反映出多年生牧草育种与籽粒作物育种的一项区别：每次去叶既是当期的利用，也是下一轮生长的起点，品种因而需要在反复利用中持续形成可收获的冠层。

在草坪绿化与地表植被恢复中，评价重点则由可收获的茎叶转向草层覆盖及其稳定性。运动场草坪受到踩踏和磨损后，即使未受损区域生长旺盛，局部裸露仍会破坏草面的连续性。Głąb 等对冷季型草坪草的研究，将多年生黑麦草等草种的耐磨损表现与植株性状联系起来，体现了草层在反复机械扰动下维持覆盖的利用要求\papercite{glab2024turfwear}。在坡面防护中，Gu 等通过模拟降雨考察了多年生黑麦草与狗牙根混播植被的抗侵蚀表现\papercite{gu2025grasserosion}。这些用途赋予冠层覆盖以不同的解释：对草坪，裸露区域的大小和分布本身就是评价对象；对坡面植被，覆盖状况可用于描述地表保护，但防护效果仍需通过侵蚀观测检验；对饲用材料，冠层展开则是生长表现之一，还需结合株高与收获量判断形成了多少饲草。因此，同样是测量植株占据的地面范围，其育种意义取决于材料最终如何利用。

黑麦草的株丛结构决定了哪些生长差异可以从图像中直接测得。多年生黑麦草由多个分蘖组成，每个分蘖上着生处于不同发育阶段的叶片；分蘖排列、叶片长度与伸展方向共同决定株丛的外观。俯视时，一份材料可能由少数长叶向外伸展，另一份材料则由较多短叶形成紧凑株丛，两者占据的水平范围可以接近，冠层内部的组织数量和排列却不相同。尤其在叶片相互重叠时，新形成的组织能够增加生物量，却不一定扩大可见投影。Gebremedhin 等在多年生黑麦草育种田间试验中，结合单株区域的植被指数与株高估计生物量，正是因为水平影像和垂直结构提供了不同的信息\papercite{gebremedhin2020ryegrassbiomass}。由俯视图像提取的冠层投影范围本身具有明确的形态学含义，但若要将其用于估计干物质产量，仍需以实际收获量建立和检验二者的关系。

冠层范围看似容易界定，在黑麦草图像上却有一个具体难点：株丛最外侧的叶尖所围成的轮廓，并不等于植株组织实际覆盖的区域。舒展的株丛中，相邻叶片之间往往露出大片土壤或底布；紧凑株丛的中心则有多层叶片重叠，外围仅有少数窄叶伸出。若简单连接外侧叶尖，叶间空隙就会被计入植株面积；若只保留连续而明显的冠层主体，孤立于外围的细叶又容易遗漏。这两类误差还会随株型改变：对叶片舒展的材料，填补空隙可能明显高估覆盖；对窄叶比例较高的材料，遗漏外围叶片则可能低估投影面积。因此，黑麦草冠层测量既要保留伸出株丛的可见叶片，也要保留叶片之间真实存在的背景。用于不同材料比较时，还应明确所报告的是视野内的覆盖比例，还是经尺度标定后的物理投影面积。

刈割进一步改变了这些结构性状的时间含义。刈后初期，已有分蘖上的叶片伸长即可使冠层重新展开；在较长时期内，分蘖的发生与存活也会改变株丛组成。两份材料在刈后达到相近的覆盖度，一份可能依靠较多分蘖共同形成叶层，另一份则可能主要依靠少数叶片快速伸长。Tozer 等在多年生黑麦草去叶试验中同时考察地上部生长、根系生长和水溶性碳水化合物储备，表明评价去叶后的植株表现需要结合处理方式和观测时期\papercite{tozer2024defoliation}。Johansen 等对多年生黑麦草农艺性状与基因型和环境互作的研究，也说明材料表现需放在具体环境中比较\papercite{johansen2025ryegrassgxe}。就图像能支持的判断而言，刈后某日的投影面积描述该日已恢复的覆盖范围；比较恢复快慢，则至少需要记录刈割时间、留茬条件及前后时点的可比观测。单幅图像可以测状态，恢复速率还取决于两次观测之间发生了多少变化。

除了冠层的空间范围，叶片维持绿色或转为枯黄的程度，也是多年生黑麦草材料评价中的重要性状。对饲用材料，叶片本身是被采食或收获的组织；对草坪材料，枯黄组织增加会直接改变草层色泽与绿色覆盖。Yang 等在多年生黑麦草持绿材料中分别测定叶片衰老、饲用品质及高温下的植株表现，发现所研究材料的持绿表型伴随饲用品质和耐热性的改善\papercite{yang2025ryegrassstaygreen}。这一结果之所以具有育种意义，在于叶色以外的性状也得到了独立测定。相似地，Tang 等在比较 56 份多年生黑麦草种质的盐胁迫响应时，同时考察生长、叶片水分状况、叶绿素及离子含量等指标\papercite{tang2013ryegrasssalinity}。对大规模种质调查而言，枯叶率的优势在于能够以统一尺度记录特定条件和时期的可见枯黄程度；至于材料是否同时维持生长或具有较强的离子调节能力，则不能仅依据外观比例判定。

图像枯叶率的测量难点，在于分子与分母都要从颜色和形态复杂的同一株丛中确定。人工目测根据整株外观判断枯黄程度，俯视图像所记录的则是枯黄组织占可见冠层投影的比例。下层枯叶被上层绿叶遮挡时，它不会出现在当前视角的统计中；枯叶伸向株丛外围、与裸土或残茬相邻时，又可能因外观相近而被划入背景。Moreno 等利用多年生黑麦草与梯牧草的 RGB 图像区分活体植被、枯死材料和裸土，发现开花初期的枯死材料与裸土覆盖尤其难以准确估计\papercite{moreno2025forageseg}。这对黑麦草枯叶率测量有直接启示：如果先按绿色特征确定植株范围，部分枯黄叶片在组织分类之前便已丢失；如果把与枯叶颜色接近的土壤纳入冠层，又可能高估枯黄比例。因此，可靠的测量需要同时满足两个条件：冠层边界应涵盖可见的绿色与枯黄组织，组织状态的判别则应始终在同一植株范围内进行。比较图像值与人工读数时，也应保留两者在观测视角和判读依据上的区别。

相对于冠层范围与组织比例，叶尖数描述的是离散的可见结构。叶尖是单片叶的末端，记录其二维坐标可以同时保留数量和位置。黑麦草外围叶片的末端通常可以沿叶缘追溯，而在稠密株丛内部，一片叶被另一片叶遮挡时，其可见轮廓也会突然终止，局部外观可能与真实叶尖相似。判断该处是否为叶尖，需要参考周围叶片的走向；确定末端具体位于何处，则需要保留窄叶尖端的局部细节。刈割样本还存在另一类歧义：叶片切口同样是可见叶段的末端，却与自然形成的叶尖具有不同含义，是否计入应在标注前明确。因而，叶尖计数适宜表述为给定视角和标注规则下的可见叶尖数。保留各点位置后，才能核查模型是否漏掉真实末端、是否将交叠处误认为叶尖；仅凭图像总数接近，无法区分这两种错误。

分蘖数与叶尖数都可以表现为计数值，其生物学单位及所需的取像证据却不同。一个分蘖可以连续产生多片叶，辨认一个叶尖只能确认一片可见叶的末端，不能由此确定它属于哪一分蘖。比较独立种植材料的株丛组成，通常关心每株或每丛具有多少分蘖；描述密植草层的疏密，则可采用单位面积分蘖密度，后者还受草层内植株密度和空间分布影响。因而，选取适于观察株丛基部的视角、明确相邻枝条的辨识依据，比直接将冠层叶尖数换算为分蘖数更符合这一性状的测量要求。由此可见，黑麦草图像表型的关键不在于从一幅图像中尽可能多地给出指标，而在于使每个指标对应可辨识的植物结构：冠层投影需要保留细叶与叶间空隙，枯叶率需要在完整可见冠层内区分组织状态，叶尖和分蘖计数则分别依赖叶片末端与株丛基部的观测证据。这些差异构成了评价黑麦草表型测定方法是否适用的具体依据。


\subsection{黑麦草表型组学的研究进展与缺口}
\label{sec:ryegrass-phenomics-gaps}

多年生黑麦草的图像表型研究已涉及单株生长监测、生物量估计和混播草层组成分析。Lootens 等连续两个生长季对间隔种植的多年生黑麦草进行田间成像，将刈后株丛基部面积与人工分蘖数进行比较，并结合叶片生长测量评价再生表现\papercite{lootens2016ryegrassregrowth}。Gebremedhin 等则在育种田间试验中提取单株区域的植被指数和株高，用于估计干物质产量\papercite{gebremedhin2020ryegrassbiomass}。这些工作为利用影像比较种质生长差异提供了依据，但不同测量任务所需的参照并不相同。收获量可以检验生物量估计是否准确，却不能指出植株掩膜漏掉了哪条叶片；人工分蘖数可以用来分析株丛基部面积与分蘖数量的关系，却不能作为图像中叶尖位置的参照。因此，在评价黑麦草图像表型研究的进展时，需要区分图像中直接识别的对象、据此计算的性状，以及用于检验该性状的参照数据。

在冠层分割方面，深度学习已被用于包含多年生黑麦草的复杂草层。Bateman 等提出局部上下文卷积网络 LC-Net，对多年生黑麦草与白车轴草混播草层进行语义分割，在其测试集上获得了 81.3\% 的平均交并比\papercite{bateman2020mixedsward}。Pranga 等结合无人机 RGB 和多光谱影像，采用机器学习方法估计禾草与三叶草混播草层的物种组成\papercite{pranga2024swardcomposition}。这类研究关心的是草层中有哪些组分、各组分占多大比例。转向间隔种植的黑麦草种质时，测量对象变为一株植物，除了区分植株与背景，还需要保留株丛外围的窄叶以及叶片之间真实存在的空隙。遗漏一条外伸叶片与填补一处叶间空隙，可能使预测面积的增减相互抵消，但所得掩膜已不能准确描述该株材料。混播草层上的分割指标反映了模型在原任务中的性能，尚不足以判断它对不同株型黑麦草的单株边界是否具有相同的识别精度。

与此相关的数据不足，具体体现在图像已有的标签能否支持所关心的算法任务。以公开发布的爱尔兰禾草--三叶草数据集为例，其近地成像部分包含相机和手机拍摄的样方图像，部分样方还提供收获干物质量、禾草与三叶草等组分的质量百分比及刈后草层高度\papercite{vistamilk2024grassclover}。这些记录为估计牧草产量与组成提供了实测参照，却不能指出哪些像素属于某一株黑麦草。即使收获物中禾草的质量占比与影像中禾草的投影占比都用百分数表示，二者的分子和分母也不相同。对于单株冠层分割，更直接的参照是与原图逐像素对应的植株掩膜，并且对细窄叶缘、失绿组织和叶间背景采用一致的标注规则。在现有黑麦草图像资源中，能够同时覆盖不同株型、组织状态与成像背景，并支持上述边界评价的专用标注仍较有限。这也是已有牧草图像数量增加，而特定分割任务仍可能缺少训练与评价依据的原因。

在枯叶率测量方面，已有研究已开始区分牧草冠层内不同状态的组织。Moreno 等利用多年生黑麦草和梯牧草 RGB 图像识别活体植被、枯死材料和裸土，发现在开花初期，枯死材料与裸土的覆盖情况尤其难以准确估计\papercite{moreno2025forageseg}。对单株枯叶率而言，除组织颜色外，还要判断枯黄组织是否属于目标植株。附着于株丛外围的枯叶和散落在株间的残茬可能颜色接近，前者应计入该株的枯叶面积，后者则应排除。若前景提取主要依赖绿色特征，部分枯叶会同时从分子和分母中被剔除；若将背景残茬误计为该株枯叶，则可能使估计值偏高。因此，已选像素上的绿叶与枯叶分类性能，只能说明分类器对这些样本的判别能力。要判断完整流程能否用于种质比较，还需考察冠层范围是否正确，以及两种测量在同一份材料上的数值相差多少。人工目测与图像测量之间的相关性可以描述二者是否同向变化，但不能单独说明差值的大小，也不能把目测值当作没有误差的真实值。

可见叶尖定位与计数面临的是另一类标注不足。水稻和小麦已有可用于训练与比较的叶尖点标注，相关研究既利用三维植物模拟生成小麦苗期的叶尖标注，也采用点预测模型从田间图像提取叶尖密度\papercite{han2025fomo4wheat}\papercite{li2023selfsupervisedphenotyping}\papercite{yang2025cerealdensity}。这些数据为识别细长叶片的末端提供了基础，但模型在粮食作物图像上的计数误差与定位精度，并不能直接代表其在多年生黑麦草株丛中的表现。黑麦草多片叶常从相近位置向外伸展，一片叶在交叠处被遮挡后，其可见叶段也会突然终止，局部轮廓与自然叶尖相似；在刈后图像中，还需按统一规则区分叶片切口与自然末端。对这类样本，只有每幅图像的总数而没有叶尖坐标，便难以发现一个真实叶尖被漏掉、另一处交叠边缘被误认所造成的误差抵消。在已有黑麦草研究中，针对单株可见叶尖、以统一点标注同时检验总数与位置的公开比较仍较有限。因而，目前证据薄弱的既包括模型在黑麦草图像上能否估准叶尖数，也包括其预测点能否准确落在相应的可见末端。

单株分蘖计数需要与可见叶尖定位与计数分开讨论。叶尖是叶片的末端，分蘖则是构成株丛的枝条；标出一枚叶尖并不能确定它属于哪个分蘖。多年生黑麦草研究已有人工分蘖调查，也已有借助成像提高调查效率的尝试。Lootens 等以人工分蘖数检验其与刈后株丛基部面积的关系，所评价的是面积作为株丛扩展指标的适用性，并不要求算法在图像中逐一定位分蘖\papercite{lootens2016ryegrassregrowth}。Käch 等报道的 Tillerbox 通过标准化取像评价草地样方的分蘖密度，其统计单位是单位地表面积，而单株种质调查所要确定的是一株或一丛具有多少分蘖\papercite{kach2024tillerbox}。在其他禾本科植物中，TillerPET 利用收获后暴露的水稻稻茬测量分蘖数与株丛紧凑度，Ferm 等也已将深度学习用于翦股颖属种间杂交材料的分蘖计数\papercite{zhou2025tillerpet}\papercite{ferm2026bentgrasstillers}。这些成果为选择取像部位和计数方法提供了参照，但要评价黑麦草单株分蘖数，仍需在能够辨识株丛基部枝条的图像上，将识别结果与同一株丛的人工逐蘖计数对应起来。样方分蘖密度、株丛基部面积和冠层可见叶尖数均不能替代这一参照。

综上，现有研究已经证明图像方法能够服务于多年生黑麦草的生长监测与草层评价；对于以单株种质为单位的精细表型测量，尚待解决的问题主要集中在专用标注与任务验证上。冠层分割需要判明模型对细叶、枯叶及叶间背景的误差是否随株型改变；枯叶率测量需要区分冠层范围误差与组织判别误差，并考察二者对种质测量值的共同影响；可见叶尖定位与计数则需要在黑麦草点标注上分别检验数量和位置。相近禾本科作物的公开数据为方法比较提供了起点，而上述性状能否在黑麦草材料上得到稳定、可复核的测量，仍取决于目标作物自身的标注与验证依据。

\subsection{研究的目的与意义}
\label{sec:research-purpose-significance}

针对多年生黑麦草单株冠层与可见叶尖的专用标注不足，以及现有视觉方法在细长叶片交错的株丛中难以兼顾冠层边界提取与叶尖精确辨识的问题，本研究以多年生黑麦草种质为对象，开发基于深度学习的冠层分割与可见叶尖定位方法，并分别用于图像枯叶率测量和叶尖计数。两项任务取自同一类植株图像，但需要保留的证据不同：测定枯叶率，首先要确定哪些像素属于目标植株，再区分其中的绿色与枯黄组织；统计可见叶尖，则要判断叶片轮廓在何处真正终止，并记录各末端的位置。水稻、小麦与黑麦草在细长叶形和叶片遮挡方面的相似性，为方法开发提供了具有共同识别难点的公开样本。本研究先在公开作物数据上比较算法性能，再结合黑麦草图像构建与两项任务对应的标注数据集，考察模型在目标草种上的适用性。

在冠层分割方面，本研究提出双向多尺度语义分割网络（Bidirectional Multi-scale Segmentation Network，BiMSeg），旨在识别冠层主体的同时，保留向株丛外围伸展的细窄叶片，并区分相邻叶片之间裸露的背景。经过公开数据集上的分割性能比较，进一步以人工绘制的黑麦草掩膜为初始样本，结合模型预标注与人工校订，构建黑麦草冠层分割数据集 RyeSEG。标注时将可辨识的绿叶与枯叶均纳入植株前景，使前景的定义取决于组织是否属于目标植株，而不取决于其是否仍然保持绿色。这一区分对枯叶率测量尤为重要：若在提取冠层时遗漏枯黄叶片，即使后续颜色分类完全正确，被遗漏的组织也无法计入枯叶率。已有牧草图像研究中观察到的枯死组织与裸土混淆，进一步说明了先界定植株范围的必要性\papercite{moreno2025forageseg}。

在此基础上，本研究采用极端梯度提升算法（Extreme Gradient Boosting，XGBoost）\papercite{chen2016xgboost}对冠层内的像素进行绿叶与枯叶分类，以枯叶像素数占可见冠层前景像素总数的百分比表示图像枯叶率。为考察这一测量能否反映种质间的枯黄程度差异，在盐胁迫和高温试验中，分别按种质配对比较图像枯叶率与人工目测枯叶率，同时考察材料排序的接近程度与逐份材料的数值差异。在每项试验内，使用相同的种质和基因型数据，分别以两种枯叶率开展全基因组关联分析，以比较不同测量方式所得表型的关联结果。由此，图像枯叶率既可返回原图核查，也可作为数量性状用于群体分析；算法的测量表现与表型的遗传关联结果分别评价，两者各有相应的参照。

在可见叶尖定位与计数方面，本研究提出上下文引导的叶尖定位与计数网络（Context-Guided Network for Leaf-Tip Localization and Counting，CGTipNet）。针对叶片在交叠处形成的轮廓中断与自然叶尖局部外观相似的问题，模型利用周围叶片的走向辅助辨识末端，并保留高分辨率特征中的局部细节以确定叶尖坐标。模型以同一组预测点给出位置和数量，使计数误差较小的图像仍可进一步检查预测点是否落在真实叶尖附近。在公开水稻和小麦叶尖数据上比较定位与计数性能后，以人工点标注的黑麦草图像为基础，结合模型辅助标注构建黑麦草叶尖点标注数据集，为模型向黑麦草株丛图像迁移提供目标样本和人工参照。这里测定的是给定视角下可辨识的叶片末端数量。点定位思路也可为后续单株分蘖计数提供方法基础，但分蘖的辨识部位在株丛基部，其适用性须通过独立的分蘖标注与人工计数另行检验。

本研究的意义在于，使黑麦草种质调查中的两类读数都能对应到图像中可核查的对象：枯叶率对应目标植株范围内被判为枯叶的区域，可见叶尖数对应逐一标出的叶片末端。当材料间出现数值差异时，可以回到原图判断差异来自组织状态、株丛形态，还是来自边界误分或叶尖漏检。这为大规模种质材料在统一口径下开展表型测量、复核异常记录和比较胁迫表现提供了依据。对于高羊茅、狗牙根、早熟禾等同样具有狭长叶片或密集草层的禾草，黑麦草研究所形成的冠层标注、局部特征辨识及性状验证方法，也可为其高通量表型提取提供方法基础。